\section{A positive comparison for a region without a character}\label{cs:probe}

Consider a connected region of product junctions with no distinguished
original character at a junction. Each internal edge carries a turn.
Choose one outgoing edge. Every junction then has at least two incoming
branches. Its boundary inputs are \(n\ge2\) polynomials \(F_{i,N}\)
of type \(L\). The resulting expression, before its outgoing turn,
is denoted by \(A(F_1,\ldots,F_n)\), and its outgoing operator by
\(Q_N=\Turn A(F_1,\ldots,F_n)\).

The finite expression uses only products and turns; every input occurs
once. It can be constructed from two-input products by turns, products
of completed expressions and single-input attachments. Its type refers
to its pure positive states, not to the range of its signed Newton
coefficients.

\begin{definition}[The comparison arrays]
Let \(c_i\in\mathbb R\) be fixed constants. The input arrays satisfy
the comparison hypothesis if a common \(e_N\to0\) gives
\begin{equation}\label{cs:cone}
 |\coef{r}{F_{i,N}}|\le e^{N(c_i+e_N)}w_{N,r},
       \qquad 0\le r\le d,\quad 1\le i\le n.
\end{equation}
Replace every input coefficient by its absolute value, and apply the
same products and turns; call the resulting operator \(A_N^{\mathrm{abs}}\).
It is a positive coefficient comparison, not the pointwise absolute
value of the polynomial \(A_N\).
\end{definition}

Introduce the auxiliary polynomial
\begin{equation}\label{cs:probe-polynomial}
 W_N=\sum_{r=0}^{d}w_{N,r}e_r.
\end{equation}
It serves only in a norm inequality. It does not add a crossing circle
to the link. A \emph{band sector} chooses, for each input and for this
comparison polynomial, either the low indices \(r\le M\) or the high
indices \(r>M\). These bands are not the physical core and alias.

Let \(A^0\) be the same expression with every input replaced by the pure
complete leaf \(h_N=\Turn L_N\). Write \(\Phi_A\) for the complete
action of the closed reference expression \(\phi(h_N A^0)\), and set
\begin{equation}\label{cs:Omega}
 \Omega=\sum_i c_i+\Phi_A.
\end{equation}
The reference has \(n+1\ge3\) effective leaves. It has no flagged
junction and is strict; these facts and its stationary boundary indices
are supplied by \cite{Profiles}.

\begin{theorem}[Positive comparison with shared indices]\label{cs:unflagged}
For the fixed component and inputs \eqref{cs:cone},
\begin{equation}\label{cs:positive-bound}
 \normJ{Q_N}\le CN^A\phi(W_N A_N^{\mathrm{abs}})
                      \le e^{N(\Omega+o(1))}.
\end{equation}
Every band sector with a high input or a high comparison index has
upper bound \(e^{N(\Omega-\gamma)}\), for some \(\gamma>0\) and all
sufficiently large \(N\). The constants are uniform over the specified
input families.

The all-low comparison is concentrated at the unique strict maximum
of the original reference action. Its effective input indices and
comparison index are in \((1/4,1/2)\). Outside any fixed neighborhood
of that maximum it has a fixed exponential loss.
\end{theorem}

\begin{proof}
Positivity of \eqref{cs:product} and exact coefficient reversal give
\(|\coef{r}{A_N}|\le\coef{r}{A_N^{\mathrm{abs}}}\).
By the full pairing,
\(\coef{r}{Q_N}=\phi(e_r A_N)\). Hence
\begin{equation}\label{cs:positive-mass}
 \sum_r w_{N,r}|\coef{r}{Q_N}|
      \le\phi(W_N A_N^{\mathrm{abs}}).
\end{equation}
This proves the first inequality in \eqref{cs:positive-bound} by
\eqref{cs:basis-norm}. It holds separately in each specified band
sector and retains the first derivatives of all high basis elements.

It remains to bound the positive scalar. For low input and comparison
indices its upper leaf profile is \(g\); for high indices it is zero.
The common additive constants are \(\sum_i c_i\). Set
\[
 \widetilde g(x)=g(\max(x,1/4)),\quad 0\le x\le1/2,
 \qquad k(x)=-2J(x/2),\quad 0\le x\le1.
\]
On these domains \(g\le\widetilde g\) and \(0\le k\).
Thus the scalar supporting-plane estimates of
Lemmas~\ref{cs:two-leaf-gap} and~\ref{cs:attachment-gap} apply to the profiles of
\eqref{cs:cone}. We give the propagation argument to specify why
their applicability does not depend on the origin of the arrays.

First suppose \(n+1\ge4\). The unflagged closed junction tree can be
rooted between two completed expressions. If there is an internal
edge, each side has at least two leaves, because every internal vertex
has valence at least three. If there is only one junction, split its
at least four leaves into two products with at least two factors each.
Products, associativity and \eqref{cs:adjoint} realize this change of
root as an exact identity of full colored tensors. Every external band
stays attached to its original array. No tangle isotopy is asserted.

At a two-leaf seed with selected output \(b\), write its pure profile
as \(f(b)\), with slope \(\lambda=f'(b)\). The effective boundary
parameters have \(1/2<X<1\). The two-leaf scalar inequalities give
\begin{equation}\label{cs:seed-affine}
 \text{nonempty high profile at output }z
       \le f(b)+\lambda(z-b)-1/100.
\end{equation}
The all-low case has the same tangent without a loss. At a turn the
output index changes from \(r\) to \(d-r\), so the tangent slope
changes sign. At a product of two completed children, add their affine
bounds to the supporting plane of the concave function \(m\).
Stationarity cancels both child slopes. The outgoing slope is the
derivative of the parent profile, and inherited losses add.

For a single-input attachment, use the child's entire affine bound,
not an elementary-leaf domain for that child. If the new input is low,
the tangent of \(\widetilde g\) preserves the child's loss. If it is
high, the single-attachment inequality supplies a further loss
\(1/2500\). Its hypotheses are the effective stationary inequalities
\(1/2<X<1\) and \(0<Y,Y_o<1\) from the pure completed construction.
These hypotheses concern the fixed reference profiles. They impose no
additional condition on the coefficient arrays being bounded.

At closure the two reference outputs sum to one and their slopes
agree. The exact finite indices sum to \(d/N=1-1/N\), so the
remaining affine term is \(-\lambda/N\). After multiplication by
\(N\), it is a bounded constant. Every nonempty high assignment
has inherited at least one of the positive losses above. This proves
the required strict rate bound for \(n+1\ge4\).

For \(n=2\), the closed reference has three effective leaves.
There is a shorter direct argument. Factorial reflection
\(S_rS_{d-r}=N\) in \eqref{cs:product} gives the symmetric tensor
\begin{equation}\label{cs:trivalent}
 \phi(e_i e_j e_k)=
 \frac{S_{i+j}S_{i+k}S_{j+k}}
 {N S_iS_jS_kS_{i+j+k-d}}
\end{equation}
when
\(i+j,i+k,j+k\le d\) and \(i+j+k\ge d\), and zero otherwise.
In particular two high indices force the coefficient to vanish:
their sum is at least \(2M+2>d\), at either parity. For exactly
one high index, group it with a low index, leaving the other low
index for closure. The pure selected point is
\[
 a=\frac{\arctan2}{\pi},\quad b=1-a,\quad X=1/2,\quad
 \lambda=g'(a)=-\log4.
\]
The endpoint case of the two-leaf inequality gives
\(f(b)+\lambda(z-b)-1/100\). The remaining low profile has tangent
\(g(a)+\lambda(u-a)\). Since \(z+u=d/N\), their sum is at most
\(\Phi_A-1/100+\log4/N\). This proves the high-sector claim also
in the three-leaf reference. No estimate for three physical aliases
is needed here: these are coefficient bands, and the two-high
vanishing is an exact polynomial identity.

For the all-low sector, the uniform closed upper action is
\begin{equation}\label{cs:reference-action}
 \sum_i c_i+g(u)+\sum_i g(a_i)+\sum_v m(x_v,y_v,z_v),
\end{equation}
with the reference's original shared indices. Its unique strict
maximum is \(\Omega\). The stationary input and comparison indices
lie in \((1/4,1/2)\), including the value \(a\) just displayed in
the three-leaf case. These are the strict optimizer and effective
boundary-range results of \cite{Profiles}. Compactness gives a strict
loss outside each fixed neighborhood of that maximum.

There are only polynomially many finite states, because the component
is fixed. The factorial rates are uniform on the whole closed support.
Any limit of finite states satisfies the limiting turn and closing
constraints, so continuity of the actions transfers the compact gaps
to the displaced finite grids. Powers of \(N\), the common error in
\eqref{cs:cone} and the fixed number of sectors are absorbed by
reducing the loss. This proves all the assertions.
\end{proof}

\begin{corollary}[Transfer of an input coefficient loss]\label{cs:loss-transfer}
Let \(a_i^*\) be the selected index of input \(i\) in
\eqref{cs:reference-action}. Suppose that, on a fixed interval \(I_i\)
containing \(a_i^*\) in its interior, a specified subfamily improves
\eqref{cs:cone} to
\begin{equation}\label{cs:input-loss}
 |\coef{r}{F_{i,N}}|\le e^{N(c_i+g(r/N)-\eta+o(1))},
 \qquad r/N\in I_i,\quad \eta>0.
\end{equation}
Then \(\normJ{Q_N}\le e^{N(\Omega-\eta')}\) for some \(\eta'>0\)
on that subfamily, for every sufficiently large \(N\).
\end{corollary}
\begin{proof}
The high sectors already have a fixed gap. In the all-low sector,
choose a neighborhood of the strict maximum whose \(i\)-th coordinate
lies in \(I_i\). Inside it the improved coefficient removes
\(\eta\) from every state rate; outside it compactness gives a gap.
Use \eqref{cs:positive-mass} and reduce the minimum of these losses.
The arrays can be complex or identically zero; no lower bound for
them is used.
\end{proof}

\begin{proposition}[Angular localization at the reference rate]\label{cs:probe-strip}
The canonical outgoing polynomial in Theorem~\ref{cs:unflagged}
satisfies
\begin{equation}\label{cs:Q-strip}
 |Q_N(2\cos(\theta+\ii y))|
 \le e^{N(\Omega-c(\theta-\pi)^2+Cy^2+o(1))}
\end{equation}
on a fixed strip, with \(c,C>0\), uniformly in its physical inputs.
\end{proposition}
\begin{proof}
Choose a compact interval \(I\subset(0,1/2)\) around the selected
comparison index. Separate the actual state expansion of \(Q_N\)
into an all-low part with output index \(r/N\in I\), and its
complement. The weighted coefficient mass of the complementary part
has a strict gap by Theorem~\ref{cs:unflagged}. It is a canonical
polynomial, so \eqref{cs:ellipse} bounds its strip rate by
\(\Omega-\gamma+|y|\).

For the first part apply \eqref{cs:low-gaussian} to each output
\(e_r\). Its factor \(w_{N,r}\) is already present in
\eqref{cs:positive-mass}; no further exponential cost is added.
This gives \eqref{cs:Q-strip} for that part. Finally choose
\(y_0<\gamma/4\) and \(c\pi^2<\gamma/4\), reducing the angular
constant if needed. The complementary part then fits under the same
bound. The two contributions together prove the assertion.
\end{proof}
